\documentclass[10pt,a4paper]{amsart}
\usepackage[margin=19mm]{geometry}
\usepackage{amsmath,amssymb,mathtools}
\usepackage[scr=rsfs,cal=euler]{mathalfa}
\usepackage{xfrac}
\usepackage[unicode,hidelinks,bookmarks=false]{hyperref}
\newcommand{\ie}{i.e.\ }
\newcommand{\cf}{cf.\ }
\newcommand{\resp}{respectively\ }
\newcommand{\Subsection}{\subsection}
\providecommand{\nobreakdash}{\nobreak\hbox{-}}
\setlength{\emergencystretch}{2em}
\multlinegap0pt
\usepackage{amssymb}
\usepackage{xcolor}
\usepackage{algorithm}\usepackage{algpseudocode}
\usepackage{mathtools}
\newtheorem{proposition}{Proposition}
\title{Twist-structures isomorphic to modal Nelson lattices}
\author{Paula Mench\'on, Ricardo O. Rodriguez}
\date{Source: arXiv:2504.08109v2 (CC BY 4.0)}
\begin{document}
\maketitle

\noindent\emph{Source and licence.} Twist-structures isomorphic to modal Nelson lattices. Version arXiv:2504.08109v2 (CC BY 4.0), distributed by arXiv under the Creative Commons Attribution 4.0 International licence, http://creativecommons.org/licenses/by/4.0/, as recorded in the arXiv licence metadata for this exact version and shown on the abstract page https://arxiv.org/abs/2504.08109v2. Full licence notice: \texttt{licenses/arxiv-2504.08109-CC-BY-4.0.txt}.\par\smallskip

\noindent\textit{Editorial scope.} Counted unit: the article's proposition proposition (source0.tex lines 1191--1193). The article's complete original proof of it is delivered with it (source0.tex lines 1194--1204, one \texttt{proof} environment of 11 lines). No mathematics is written, repaired or re-derived: the statement and the proof are the article's own lines, in the article's own notation, and no retained line is substituted or re-flowed.  No further source-local context is needed: every cross-reference of every retained line resolves inside the two delivered spans.  Nothing was omitted from any retained span, so the package carries no omission record.  No retained line contains a citation, so the wrapper carries no bibliography and no bibliography entry.  No retained line contains a figure environment, an image-inclusion command or drawing code, so no image is delivered and the package declares no asset dependency.  Editorial formatting only, in addition: an AMS \texttt{amsart} preamble that restates the article's own declarations and macros used by the retained lines, namely the environments \texttt{proposition} on separate counters, together with the standard packages those lines need.  The retained environments print in this document as Proposition 1 in order of appearance, whereas the article prints the same items with the numbering of its own full document.  The complete native source of the article as distributed in the arXiv e-print for this exact version is bundled unchanged as \texttt{sources/176347/source0.tex}.
\begin{proposition}
    Let $f \colon X_1\to X_2$ be an ME-morphism	between two ME-spaces $\mathcal{X}_1=\langle X_1,\tau_1,\leq_1,\eta_1,\eta_2\rangle$ and $\mathcal{X}_2=\langle X_2,\tau_2,\leq_2,\eta'_1,\eta'_2\rangle$. Then, the map $h(f) \colon \mathcal{U}(X_2)\to \mathcal{U}(X_1)$ defined by $h(f)(U)=f^{-1}[U]$ is a modal Heyting algebra homomorphism. 
\end{proposition}

\begin{proof}
It is known that $h(f)$ is a Heyting algebra homomorphism. Let $U\in \mathcal{U}(X_2)$. Then,
    \begin{align*}
    \square_{\eta_1}(h(f)(U))&=\square_{\eta_1}(f^{-1}[U])\\&=\{x\in X_1:f^{-1}[U]\in\eta_1(x)\}\\&=\{x\in X_1:U\in\eta'_1(f(x))\}\\&=\{x\in X_1:f(x)\in\square_{\eta'_1}(U)\}\\&=f^{-1}[\square_{\eta'_1}(U)]\\&=h(f)(\square_{\eta'_1}(U)).
    \end{align*}
On the other hand,
    \begin{align*}
        \Diamond_{\eta_2}(h(f)(U))&=\Diamond_{\eta_2}(f^{-1}[U])\\&=\{x\in X_1:X_1\setminus f^{-1}[U]\in\eta_2(x)\}\\&=\{x\in X_1:X_2\setminus U\in\eta'_2(f(x))\}\\&=\{x\in X_1:f(x)\in\Diamond_{\eta'_2}(U)\}\\&=f^{-1}[\Diamond_{\eta'_2}(U)]\\&=h(f)(\Diamond_{\eta'_2}(U)).
    \end{align*}
    Since $U$ is arbitrary, it follows $h(f)$ is a modal Heyting algebra homomorphism.
\end{proof}

\end{document}
